% =====================================================================
\lezione{28/09/2026}{GDPR (II): legittimo interesse, diritti, responsabili, accountability, art.~32 e DPIA}
  {Prof.~Giuseppe Vaciago}
  {interessato; legittimo interesse e videosorveglianza; diritti dell'interessato; sanzioni; data
   retention; DPO; responsabile e sub-responsabile, Data Privacy Framework; accountability e stato
   dell'arte; art.~32 e misure tecniche e organizzative; risk assessment, privacy by design e by default,
   DPIA e consultazione preventiva}
\label{lez:2026-09-28}
% =====================================================================

\begin{nota}[Fonte]
Trascrizione automatica della registrazione, revisionata. Durante la pausa la registrazione contiene
conversazioni in italiano non pertinenti, ignorate. Slide di riferimento: \emph{02.~Cybersecurity Laws
and Regulations -- II Privacy}, dal legittimo interesse fino al \emph{risk assessment} (la parte sul
data breach, moduli M4--M5, è rinviata). In coda, uno scambio in italiano con uno studente collegato
da remoto sul significato di \emph{formal assessment}.
\end{nota}

\section{Sintesi}
Seconda parte del GDPR. Dopo un riepilogo e la definizione di \textbf{interessato}, il docente chiarisce
che il \textbf{legittimo interesse} è del titolare, non dell'interessato, e passa ai diritti, alle
\textbf{sanzioni} (con la distinzione 20M/4\% e 10M/2\%), alla conservazione dei dati e alla figura del
DPO. Ampio spazio al \textbf{responsabile} e al \textbf{sub-responsabile} (con un caso reale di
misurazione del piede con IA e il Data Privacy Framework UE--USA) e all'\textbf{accountability},
legata allo \textbf{stato dell'arte}. Chiudono l'art.~32, le misure tecniche e organizzative e una
prima introduzione al risk assessment (probabilità e gravità), alla privacy by design e by default e
alla DPIA con consultazione preventiva (esempio dell'app Immuni).

% ---------------------------------------------------------------------
\section{Organizzazione: i laboratori}
\begin{itemize}
  \item Funzionamento: il docente presenta un caso (il primo sarà un \textbf{data breach}), gli studenti
    lavorano in piccoli gruppi e poi presentano. I migliori gruppi (3--5, a seconda dei numeri)
    ottengono \textbf{punti extra} per l'esame.
  \item Due obiettivi: imparare un \textbf{ragionamento giuridico} semplice, diverso da quello tecnico,
    e \textbf{parlare in pubblico}.
  \item La grande maggioranza dei presenti è interessata. Decisione annunciata dopo la pausa: i
    \textbf{prossimi due laboratori} si fanno \textbf{da remoto}, poi \textbf{dal 16 ottobre in
    presenza}. Niente modalità ibrida: chi è collegato resta escluso dal lavoro di gruppo.
\end{itemize}

% ---------------------------------------------------------------------
\section{Riepilogo e definizione di interessato}
Riepilogo della lezione del 25/09: ruoli, bilanciamento, Google Spain (``importante l'esempio, non il
caso''), principi, trattamento (``ogni attività può essere un trattamento''), esenzione domestica,
quattro pilastri, tre categorie di dati, consenso.

L'\textbf{interessato} (\emph{data subject}) è la \textbf{persona fisica} a cui si riferiscono i dati,
mai una persona giuridica. Il GDPR va rispettato da chi tratta i dati (titolare, responsabile); la
persona fisica è chi il GDPR protegge. Siamo tutti interessati.

% ---------------------------------------------------------------------
\section{Il legittimo interesse}
\begin{esame}
Il legittimo interesse è \textbf{del titolare}, non dell'interessato: l'interessato ha
\emph{diritti fondamentali}, non ``legittimi interessi''. È un errore frequente e una possibile domanda
d'esame.
\end{esame}

\paragraph{Videosorveglianza.} Chiedere il consenso a ogni persona ripresa è impossibile (e in alcuni
luoghi, come stadi e concerti, le telecamere sono imposte dalla legge). La base giuridica è il
legittimo interesse, frutto di un bilanciamento tra la privacy di chi passa e l'interesse a
ricostruire, ad esempio, un incidente stradale. Le slide ricordano che il bilanciamento resta
obbligatorio anche nei casi tipici (banche, gioiellerie, distributori di benzina).

\paragraph{Marketing ai clienti.} Un negozio di abbigliamento manda offerte a chi ha già comprato: da un
lato il suo legittimo interesse verso un cliente, dall'altro il diritto di non essere disturbati. Serve
il \textbf{balancing test}, caso per caso: è ragionevole proporre prodotti simili a quelli acquistati,
non qualsiasi cosa (altrimenti è spam). Quando la base è il legittimo interesse esiste \textbf{sempre}
il \textbf{diritto di opposizione} (il link ``unsubscribe'').

\paragraph{Il legittimo interesse non vale per le categorie particolari.} Su domanda di uno studente:
per i dati dell'art.~9 non si può usare il legittimo interesse (si veda l'elenco delle basi dell'art.~9
nelle slide). Un'immagine di videosorveglianza che mostra una caratteristica visibile non è per questo
un dato sanitario; una patologia non visibile sì.

\subsection{Europa e Stati Uniti: due culture}
Negli USA si può fare quasi tutto con i dati dei propri clienti, tranne venderli; si ha invece paura
dello Stato (``Big Brother''). In Europa è il contrario: si accetta il trattamento pubblico e si teme
quello privato. Per il docente entrambe le posizioni sono sbagliate: grandi aziende tecnologiche (e ora
quelle dell'IA) hanno un peso economico superiore a molti Stati, e l'unica strada è la
\textbf{cooperazione pubblico--privato} e il bilanciamento. Esempio: i dati sanitari dei wearable
potrebbero alimentare la prevenzione, ma il quadro privacy lo rende difficilissimo, mentre molti
dispositivi a basso costo cedono comunque i dati a Paesi non democratici.

% ---------------------------------------------------------------------
\section{I diritti dell'interessato}
\begin{itemize}
  \item \textbf{Accesso}: riguarda solo i \emph{propri} dati. Esempio del docente: lo studente può
    accedere ai dati che ha fornito all'esame, ma la valutazione del docente non è un suo dato.
  \item \textbf{Rettifica} dei dati inesatti.
  \item \textbf{Cancellazione} (diritto all'oblio).
  \item \textbf{Limitazione del trattamento}: poco usata e difficile da applicare.
  \item \textbf{Portabilità}: ricevere i dati in formato strutturato. Fondamentale in teoria, molto
    complessa in pratica: alla prima richiesta di un attivista austriaco, Facebook consegnò i dati su
    migliaia di pagine stampate. Per il docente la portabilità richiede interoperabilità e open
    source.
  \item \textbf{Opposizione}: chiedere che cessi un trattamento, salvo motivi prevalenti del titolare.
\end{itemize}

% ---------------------------------------------------------------------
\section{Sanzioni}
In generale fino a \textbf{20 milioni di euro o il 4\%} del fatturato mondiale annuo; l'autorità
usa il criterio che dà l'importo più alto (per un gruppo come Google, il 4\%).

\begin{esame}
\textbf{Domanda insidiosa}: per le violazioni delle \textbf{misure di sicurezza}, quindi anche per i
\textbf{data breach}, il tetto è la metà: \textbf{10 milioni di euro o il 2\%} del fatturato. ``Avrete
probabilmente una domanda sul data breach all'esame.''
\end{esame}

% ---------------------------------------------------------------------
\section{Compliance in quattro regole e data retention}
Le quattro regole operative delle slide (non fondamentali per l'esame, ma utili): analizzare quali dati
si raccolgono e dove (\textbf{registro dei trattamenti}, seguendo il flusso dei dati); verificare se il
tempo di conservazione è giustificato (\textbf{policy di conservazione e cancellazione}); informare i
clienti su come modificare o cancellare i dati (informativa); monitorare chi accede ai dati.

\paragraph{Per quanto tempo si conserva un CV?} Di norma \textbf{da uno a due anni} per un'azienda;
circa \textbf{sei mesi} per le agenzie per il lavoro, che ne ricevono moltissimi. Perché cancellare?
\begin{itemize}
  \item \textbf{sicurezza}: in caso di fuga di dati, ciò che non serviva non doveva esserci;
  \item \textbf{esattezza}: un CV vecchio non riflette le competenze attuali e può danneggiare la
    persona.
\end{itemize}
Ogni categoria di dati ha un tempo diverso (in Italia, ad esempio, gli attestati di formazione possono
servire al lavoratore per decenni a fini previdenziali). E la retention va coordinata con i
\textbf{backup}: un backup che conserva dati oltre i termini viola la norma. Per il docente la
retention è ``il vero caos del GDPR'' ed è un problema tecnico, quindi degli ingegneri.

% ---------------------------------------------------------------------
\section{I soggetti del GDPR e il DPO}
Interessato, titolare, responsabile, \textbf{DPO} (\emph{Data Protection Officer}), autorità di
controllo. Il DPO (art.~37) può essere interno o esterno ed è \textbf{obbligatorio} per:
\begin{itemize}
  \item autorità e organismi pubblici;
  \item titolari o responsabili le cui attività principali richiedono il \textbf{monitoraggio regolare
    e sistematico} degli interessati su larga scala;
  \item titolari o responsabili le cui attività principali consistono nel trattamento \textbf{su larga
    scala} di categorie particolari o di dati giudiziari.
\end{itemize}
Per il docente, nel 2026, con NIS2, DORA, Cybersecurity Act, AI Act e la \emph{shadow AI}, il DPO
dovrebbe esserci in ogni azienda. Servono competenze giuridiche ma anche \textbf{tecniche} (reti,
storage, flussi di dati): oggi il 90\% dei DPO sono avvocati, ma il ruolo sarebbe adatto agli
ingegneri. Il DPO deve essere \textbf{indipendente}: risorse adeguate, nessuna istruzione dal titolare
sui propri compiti, nessuna penalizzazione per averli svolti, nessun conflitto di interessi; non è
personalmente responsabile della non conformità. Nella pratica subisce le pressioni del CEO, che vede
compliance e sicurezza solo come costi. Il docente cita una truffa da decine di milioni di euro ai
danni di Fideuram con la voce di un avvocato clonata dall'IA: episodi così stanno cambiando la
percezione. La decisione finale resta comunque del titolare.

% ---------------------------------------------------------------------
\section{Responsabile e sub-responsabile}
\paragraph{Responsabile (art.~28).} Il criterio per distinguerlo dal titolare è il \textbf{potere
decisionale} sul trattamento. Il responsabile deve offrire \textbf{garanzie sufficienti} su misure
tecniche e organizzative adeguate (art.~28(1)). Il contratto (il \emph{Data Processing Agreement},
art.~28(3)) deve indicare: oggetto e durata, natura e finalità, tipo di dati e categorie di
interessati, obblighi e diritti del titolare.

\subsection{Caso: la misura del piede con l'IA}
Un'azienda francese offre ai siti di calzature un servizio: si fotografa il piede con lo smartphone e
l'IA indica la taglia giusta (riducendo i resi). L'azienda sostiene di essere \textbf{titolare}, perché
vuole usare le foto per migliorare il servizio e addestrare l'IA.
\begin{itemize}
  \item Ragionamento del docente: bisogna guardare all'\textbf{interessato}. Il cliente vuole comprare
    un paio di scarpe, non misurarsi il piede; la finalità è del marchio di calzature, che è il
    titolare. Il fornitore è \textbf{responsabile}.
  \item La \textbf{contitolarità} è possibile, ma nel bilanciamento non è nell'interesse del
    consumatore cedere a un terzo un database di milioni di immagini del corpo (i termini prevedono
    l'uso per l'addestramento dell'IA; esistono anche rischi legati a usi feticistici). Con un semplice
    rapporto titolare--responsabile, il titolare può vietare usi ulteriori e fissare i tempi di
    conservazione.
  \item Una decisione giuridica sui ruoli può cambiare la sorte dei dati di milioni di persone.
\end{itemize}

\paragraph{Sub-responsabile.} Il responsabile può ricorrere a un altro responsabile solo con
l'autorizzazione scritta, specifica o generale, del titolare (art.~28(2)), imponendogli gli stessi
obblighi; il responsabile iniziale resta pienamente responsabile verso il titolare (art.~28(4)).
Esempi: il fornitore francese che usa un provider di IA; un'agenzia di selezione che fa il backup dei
CV su Microsoft. Il titolare deve sapere quali sub-responsabili ci sono e \textbf{dove risiedono i
dati}: la catena (responsabile, sub-responsabile, sub-sub-responsabile\ldots) va seguita, e secondo il
docente moltissimi incidenti nascono da sub-responsabili di cui il titolare non era consapevole.

\paragraph{Data Privacy Framework.} Elenco delle aziende statunitensi aderenti al quadro UE--USA per il
trasferimento dei dati. Se il fornitore è nell'elenco con stato attivo (es.\ Microsoft, OpenAI), il
trasferimento è più semplice; se non c'è, servono accordi specifici come le \textbf{clausole
contrattuali standard}. Poiché la grande maggioranza dei fornitori di storage e trattamento è
statunitense, la verifica conta molto. Non è una domanda d'esame.

% ---------------------------------------------------------------------
\section{Accountability e stato dell'arte}
\textbf{Accountability} significa saper sempre spiegare, in modo legale e adeguato, \textbf{perché} si
è presa una decisione. Nella zona grigia bisogna poter dimostrare di stare dalla parte bianca.
L'ingegnere tende a risolvere il problema; il giurista guarda anche al processo e alle conseguenze
(esempio del docente: una ristrutturazione ``efficiente'' delle ferrovie con prepensionamenti di massa,
che decenni dopo pesa sul sistema pensionistico). Con l'IA, che decide in modo probabilistico e poco
spiegabile, tracciare il processo decisionale è ancora più importante.

\begin{esame}
Ciò che il docente vuole vedere dagli studenti: in un incidente si decide in pochi minuti e la
decisione può rivelarsi sbagliata, ma bisogna \textbf{documentarla e renderla tracciabile}, scrivendo
perché la si prende. ``Non diventate avvocati: serve risolvere il problema, ma documentate le
decisioni.''
\end{esame}

\paragraph{Stato dell'arte.} Con un budget illimitato si proteggerebbe tutto: le misure vanno valutate
con un budget limitato e devono essere \textbf{economicamente sostenibili}. Lo stato dell'arte cambia:
quindici anni fa l'\textbf{autenticazione a più fattori} era un'idea innovativa, oggi è di fatto un
obbligo per il titolare, e senza MFA si rischia la sanzione. L'interessato, come privato, è libero di
non usarla (anche se è fortemente consigliata).

\begin{esame}
Due parole da ricordare: \textbf{accountability} e \textbf{stato dell'arte} (spiegati bene
nell'art.~32).
\end{esame}

% ---------------------------------------------------------------------
\section{Art.~32: sicurezza del trattamento}
Tenendo conto dello stato dell'arte, dei costi di attuazione, della natura, dell'ambito, del contesto e
delle finalità del trattamento e dei rischi per i diritti e le libertà delle persone, titolare e
responsabile mettono in atto \textbf{misure tecniche e organizzative adeguate al rischio}. L'elenco è
\textbf{aperto} (pseudonimizzazione, cifratura, riservatezza e integrità, disponibilità e resilienza,
test periodici): prima c'era una lista chiusa di misure, ma con nuove vulnerabilità ogni settimana una
lista non regge. Bisogna decidere ogni volta e saper spiegare perché la soluzione scelta è adeguata.

\begin{integrazione}[la vecchia lista di misure]
La ``lista'' a cui il docente si riferisce è l'Allegato~B al D.lgs.~196/2003 nella versione originaria
(\emph{Disciplinare tecnico in materia di misure minime di sicurezza}), superato con l'arrivo del GDPR e
del suo approccio basato sul rischio.
\end{integrazione}

\begin{itemize}
  \item Va verificata la sicurezza del \textbf{responsabile}, che materialmente detiene i dati: la scelta
    tra provider (grandi hyperscaler statunitensi o provider europei) è una decisione caso per caso,
    con vantaggi e svantaggi di compliance e sicurezza.
  \item Chiunque agisca sotto l'autorità di titolare o responsabile tratta i dati solo su istruzione
    (art.~32(4)): servono \textbf{istruzioni e formazione}.
  \item Per ruoli e responsabilità in caso di incidente il docente richiama la matrice \textbf{RACI}
    (chi è responsabile, chi decide, chi è consultato, chi è informato): come in un'evacuazione
    antincendio, ognuno deve sapere cosa fare.
\end{itemize}

\paragraph{Misure tecniche e organizzative (report ENISA 2021).} Tecniche: cultura della
cybersicurezza, formazione, rapporti con terze parti (anche con ENISA), procedura di data breach,
accesso sicuro ai sistemi, sicurezza dei dispositivi, della rete, fisica, backup, cloud, siti web,
condivisione delle informazioni. Organizzative: modello organizzativo, informative, policy (retention,
data breach, valutazione dei rischi), privacy by design, consensi, documentazione (registro dei
trattamenti, DPA). Le misure organizzative sono ``pezzi di carta'', ma concreti.

\begin{esame}
Conoscere la \textbf{differenza tra misure tecniche e misure organizzative}. Gli elenchi delle slide
sono solo esempi e non vanno imparati a memoria.
\end{esame}

% ---------------------------------------------------------------------
\section{Risk assessment: prima introduzione}
Il docente presenta una visione ``più giuridica'' (\emph{anticipatory compliance}); Gianluca ne darà
una più ibrida, tecnico-giuridica, e la \textbf{definizione di riferimento di risk assessment} arriverà
con lui tra qualche settimana.
\begin{itemize}
  \item \textbf{Analisi preliminare}: mappare sistemi \emph{e} dati. Il tecnico guarda ai sistemi, ma
    la qualità dei dati cambia tutto.
  \item \textbf{Strategia sui dati} e \textbf{piano di remediation}: portare il rischio, ad esempio,
    da 8 a 3 su 10.
  \item \textbf{Minimizzazione e approccio by design}.
\end{itemize}

\paragraph{Dati finanziari: ENISA e GDPR.} Per il GDPR i dati finanziari sono \textbf{dati comuni};
nella metodologia ENISA di severity assessment (che si vedrà nel laboratorio) pesano più dei dati
semplici, perché nasce da una logica di cybersicurezza vicina al settore finanziario: l'obiettivo
finale del criminale è quasi sempre il denaro.

\paragraph{By design e by default.} \emph{By design}: valutare la privacy fin dal primo momento della
progettazione, non a prodotto finito (``a quel punto è troppo tardi''). \emph{By default}: per
impostazione predefinita si trattano solo i dati necessari; esempio dell'\textbf{opt-in} europeo
(serve il tuo sì) contro l'\textbf{opt-out} statunitense (si tratta finché non ti opponi).

\paragraph{Il modello in tre fasi delle slide.} Valutazione generale (artt.~24 e 32) e approccio by
design/by default (art.~25); valutazione formale (art.~35, rischio elevato); consultazione preventiva
(art.~36); compliance ed enforcement (art.~83). Il docente avverte che questa slide è
``discutibile'' e di non prenderla come definizione definitiva.

\paragraph{DPIA e consultazione preventiva: l'esempio Immuni.} L'app di contact tracing del periodo
COVID fu considerata ad alto rischio: dopo la DPIA il rischio residuo restava elevato, quindi serviva la
\textbf{consultazione preventiva} con il Garante, che indicò ulteriori misure di mitigazione prima del
via libera. Dopo l'approvazione il lavoro non finisce: la compliance va \textbf{monitorata di continuo},
perché ogni applicazione può generare nuovi rischi.

\paragraph{Probabilità e gravità.} Ogni valutazione del rischio si basa su due parole:
\textbf{probabilità} (\emph{likelihood}) e \textbf{gravità} (\emph{severity}). Le misure tecniche e
organizzative servono a mitigare il rischio così stimato; l'approccio è circolare (valutazione,
misure, test).

\subsection{Dopo la lezione: valutazione generale e specifica}
Nello scambio in italiano con uno studente collegato si chiarisce:
\begin{itemize}
  \item la \textbf{valutazione generale} (artt.~24 e 32) riguarda tutto l'ambiente IT e ogni misura di
    sicurezza, in modo continuo; la valutazione dell'\textbf{art.~35} (DPIA) riguarda un
    \textbf{trattamento specifico} (es.\ un nuovo smartwatch) ed è più dettagliata. Più che
    ``formale'', sarebbe meglio chiamarla ``specifica'';
  \item probabilità $\times$ gravità dà il \textbf{rischio inerente}; applicando le misure di mitigazione
    si ottiene il \textbf{rischio residuo}, cioè quello che resta dopo le misure. Se resta alto, si va dall'autorità, che può indicare altre
    misure o vietare il trattamento;
  \item con l'AI Act, invece, il livello di rischio dei sistemi ad alto rischio non si autodefinisce: è
    la norma a classificarli (es.\ IA sul posto di lavoro).
\end{itemize}

% ---------------------------------------------------------------------
\section{Domande aperte e punti da verificare}
\begin{verifica}
\begin{itemize}
  \item \textbf{Calendario dei laboratori}: dalla trascrizione, i due lab da remoto dovrebbero essere il
    2 e il 9 ottobre, poi in presenza dal 16 ottobre. Il 2/10 si è poi svolta solo una simulazione (vedi
    lezione del 02/10).
  \item \textbf{Sanzione a Google per Street View}: il docente cita una sanzione di 400 milioni di euro
    ``di ieri'' a Google Ireland. Notizia non verificata.
  \item \textbf{Data Privacy Framework}: secondo il docente un noto fornitore di IA non compare
    nell'elenco. Da verificare sul sito ufficiale, perché la situazione cambia nel tempo.
  \item \textbf{Fideuram}: l'importo della truffa con voce clonata non è affidabile nella trascrizione.
  \item \textbf{Modello in tre fasi}: il docente ha dichiarato la slide ``discutibile''; la definizione
    di risk assessment da usare all'esame è quella che darà Gianluca.
  \item \textbf{Riferimenti dell'art.~28}: paragrafi (1), (2), (3), (4) ripresi dalle slide; il docente
    non li ha letti in aula.
\end{itemize}
\end{verifica}
